\section{The transition-scale Green reduction}\label{sec:transition}

We now keep the complete physical field and choose a quantitative
consumer. This section integrates the child's B01, E01 and E04 results
\cite{GNPRSOuter,GNPRSLogTransfer}; their inherited field input is stated
explicitly below. All constants are finite, but their numerical values
and the first admissible height have not been evaluated.

\subsection{One clock, one carrier, one defect margin}
For a large positive center \(X\), set
\begin{equation}\label{eq:transition-band}
 L=L_X=\log(X/(4\pi)),\quad \lambda=c/L,\quad
 \tfrac12\le c\le1,\quad R_X=\frac{32}{\pi}\log X,
 \quad W_X=[X-R_X,X+R_X].
\end{equation}
The time \(\lambda\) and parameter \(c\) are fixed throughout every
spatial integral at this center. Write \(d\omega_X=K(x-X)\dd x\) on
\(W_X\) and \(q_X=\omega_X(W_X)\). This one-boundary measure is distinct
from the two-boundary probability measure in Section~\ref{sec:green}:
\begin{equation}\label{eq:window-mass}
 0<q_X\le\tfrac12,\qquad
 0\le\tfrac12-q_X\le\frac2\pi X^{-4}.
\end{equation}
Indeed \(K(t)\le\ee^{-\pi|t|/8}/8\), and integration of its two
exterior tails proves the estimate.

Use the upper-half-plane branches in
\begin{align}\label{eq:transition-carrier}
 L(z)&=\Log(z/(4\pi)),\nonumber\\
 \mathcal A(\lambda,z)&=\frac{\pi^{1/4}}{16}z^{7/4}
 \exp\left\{-\frac{\pi z}{8}
       +\frac{\lambda}{16}\left(L(z)^2-\frac{\pi^2}{4}\right)\right\},\nonumber\\
 \Theta(\lambda,z)&=\frac z4(L(z)-1)
                  +\frac{\pi\lambda L(z)}{16}+\frac{7\pi}{8},\qquad
 C_\pm=\tfrac12\mathcal A\ee^{\pm\ii\Theta},\qquad
 V_\lambda=H_\lambda/C_-.
\end{align}
These are explicit nonvanishing carriers on the open strip; they are
not definitions through the unknown zeros. The fixed consumer is
\begin{equation}\label{eq:transition-margin}
 m_*:=m(1/2,1)=\frac{\sin(\pi/8)}{2\cosh(\pi/8)}.
\end{equation}
By Corollary~\ref{cor:margin}, \(J_\lambda(X)<m_*\) excludes a zero
with \(1/2\le\Im z\le1\), \(|\Re z-X|\le1\), at this time.
This fixed-depth consumer is not asserted equivalent to RH.

\begin{lemma}[Harmonic normalization]\label{lem:carrier-harmonic}
The carrier has \(\calJ[C_-]=0\). Its real-boundary logarithm and that
of \(V_\lambda\) are absolutely integrable against \(K\), including
the origin and the negative half-line. Consequently, with
\(I_\lambda(X)=\int_\R K(x-X)\log|V_\lambda(x+\ii0)|\dd x\),
\begin{equation}\label{eq:normalized-green-exact}
 J_\lambda(X)-I_\lambda(X)
 =\int_\R K(x-X)\log|V_\lambda(x+8\ii)|\dd x
       -\log|V_\lambda(X+4\ii)|.
\end{equation}
\end{lemma}
\begin{proof}
The displayed analytic logarithm of \(C_-\) is harmonic in the strip.
At the boundary origin it is \(O(1+|\log|z||^2)\), with gradient
\(O((1+|\log|z||)/|z|)\). Remove an upper semicircle of radius
\(\epsilon\). The Green kernel vanishes linearly there, its normal
derivative is bounded, and the extra Green-formula terms are
\(O(\epsilon(1+\log^2\epsilon))\). On the far vertical sides the
logarithm grows at most linearly plus logarithmically squared, while
harmonic measure decays exponentially. Both contributions vanish.
The boundary modulus, using argument \(\pi\) at negative \(x\), is
\[
 |C_-(\lambda,x+\ii0)|=\frac{\pi^{1/4}}{32}|x|^{7/4}
 \exp\left\{-\frac{\pi|x|}{8}
   +\frac{\lambda}{16}\left(\log^2\frac{|x|}{4\pi}-\frac{\pi^2}{4}\right)\right\}.
\]
Its logarithm has a uniform unit-interval bound of the form in
Lemma~\ref{lem:unit-log}, on both boundaries. Subtracting it from the
complete field's logarithm proves the integrability assertion. Subtract
its zero Green functional from Theorem~\ref{thm:green} to obtain
\eqref{eq:normalized-green-exact}.
\end{proof}

\subsection{The complete field input and the outer baseline}
We use the fixed-buffer heat-density approximation recorded in
\cite{GNPRSHDA,GNPRSOuter}. It is an inherited internally assessed
input, not re-proved from saddle asymptotics here. On the fixed charts
recentered at each large real \(x\), \(0\le y\le8\), put
\[
 N_x=\left\lfloor\sqrt{x/(4\pi)}\right\rfloor,\qquad
 \nu_\lambda(z)=z/2+\pi\lambda/8,\qquad
 w_n(\lambda,z)=n^{-1/2}
   \exp\{\lambda[(\log n)^2-L(z)\log n]/4\}.
\]
The complete approximant is
\begin{equation}\label{eq:transition-provider}
 F_{N_x}(\lambda,z)=\mathcal A(\lambda,z)
       \sum_{n\le N_x}w_n(\lambda,z)
       \cos(\Theta(\lambda,z)-\nu_\lambda(z)\log n)
       +\ee^{-\pi^2\lambda/64}\mathcal R_{N_x}(z),
\end{equation}
where the retained correction is
\begin{align*}
 \mathcal R_N&=(-1)^N\mathcal G(z)a(z)^{-1/2}\mathcal C(a(z)-N),
 &a(z)&=\sqrt{z/(4\pi)},\\
 \mathcal G(z)&=[M(s_+)M(s_-)]^{1/2},
 &\mathcal C(p)&=\frac{\cos(2\pi(p^2-p)-\pi/8)}{\cos(2\pi p)}.
\end{align*}
Here \(M,s_\pm\) are as in Section~\ref{sec:RS}, with compatible
branches. The apparent singularities of \(\mathcal C\) are removable.
The input supplies fixed constants \(h_9,g_9,f_9\) such that, for
\(0\le\lambda L_x\le8\) after its fixed chart onsets,
\begin{align}\label{eq:transition-provider-bounds}
 |H_\lambda-F_{N_x}|&\le h_9x^{-3/4}(|C_+|+|C_-|),\nonumber\\
 |\mathcal G(z)|&\le g_9|\mathcal A(0,z)|/2,
 &|a(z)^{-1/2}\mathcal C(a(z)-N_x)|&\le f_9x^{-1/4}.
\end{align}
The correction bound is uniform in the fractional cutoff. The
fixed-buffer input includes its full Gaussian exterior. The charts
are recentered, never enlarged to radius \(R_X\).

\begin{theorem}[Uniform outer prime baseline]\label{thm:outer-baseline}
There are constants \(C_{\rm out},X_*<\infty\), independent of
\(X,c\) in \eqref{eq:transition-band}, such that
\begin{equation}\label{eq:outer-baseline}
 |J_\lambda(X)-I_\lambda(X)+P_{c,\lambda}(X)|
       \le C_{\rm out}/L_X\qquad(X\ge X_*),
\end{equation}
where, writing \(\Lambda_{\rm vM}\) for the von Mangoldt function,
\begin{equation}\label{eq:transition-prime}
 P_{c,\lambda}(X)=\sum_{n\ge2}\frac{\Lambda_{\rm vM}(n)}{\log n}
 \frac{n^{-5/2-c/4}}{1+n^{-4}}
 \cos\bigl((X/2+\pi\lambda/8)\log n\bigr).
\end{equation}
The omitted integers beyond \(Q\ge2\) cost at most
\(Q^{-3/2-c/4}/(3/2+c/4)\le(2/3)Q^{-3/2}\).
\end{theorem}
\begin{proof}
On \(|x-X|\le R_X\), \(y=4,8\), the eventual gates
\(R_X+8\le X/2\), \(L_X\ge16\) give
\(|L(x+\ii y)-L_X|\le2(R_X+8)/X\) and \(\lambda L_x<2\).
Dividing \eqref{eq:transition-provider} by \(C_-\), the dominant arm is
\[
 D_N(z)=\sum_{n\le N_x}n^{-s_\lambda(z)}
                         \ee^{\lambda(\log n)^2/4},\qquad
 s_\lambda(z)=\frac{1-\ii z}{2}+\frac{\lambda L(z)}4
                               -\frac{\ii\pi\lambda}{8}.
\]
The reflected carrier satisfies \(|\ee^{2\ii\Theta}|\ll x^{-y/2}\).
Since \(|w_n|\le n^{-1/2}\), summing the entire reflected population
costs \(O(x^{(1-y)/4})\). The normalized correction costs
\(O(x^{-(y+1)/4})\), and the remaining normalized error costs
\(O(x^{-3/4})\), by \eqref{eq:transition-provider-bounds}.

Let \(D_j=\sum_{n\ge2}n^{-5/2}(\log n)^j\), \(j=0,1,2\).
On these outer lines, \(\log n\le\Re L(z)/2\), so
\(|n^{-s_\lambda(z)}\ee^{\lambda(\log n)^2/4}|\le n^{-5/2}\).
Using \(\ee^u-1\le u\ee^u\) pays removal of the heat weight by
\(\lambda D_2/4\). The ordinary Dirichlet-series tail costs
\((2/3)N_x^{-3/2}=O(x^{-3/4})\). Replacing \(s_\lambda(z)\) by
\[
 s_y^0(x)=\frac{1+y}{2}+\frac c4
                   -\ii\left(\frac x2+\frac{\pi\lambda}{8}\right)
\]
costs at most \(D_1\lambda(R_X+8)/(2X)\). These uses of the
ordinary zeta series are all in \(\Re s\ge5/2\)
\cite{DLMFZetaOuter2026}. Thus a fixed \(C_F\) gives
\begin{equation}\label{eq:outer-field-comparison}
 |V_\lambda(x+\ii y)-\zeta(s_y^0(x))|\le C_F/L_X,
 \qquad y=4,8,\quad |x-X|\le R_X.
\end{equation}
No such absolutely convergent comparison has been used at \(y=0\).

The independent floor
\(|\zeta(s)|\ge1-D_0>1/2\) follows from
\(D_0\le7/(12\sqrt2)<1/2\). Once \(L_X>4C_F\), the logarithmic
error in \eqref{eq:outer-field-comparison} is at most \(4C_F/L_X\).
The unit-log bound in Lemma~\ref{lem:carrier-harmonic}, summed in
unit shells outside \(W_X\), gives on either boundary
\begin{equation}\label{eq:transition-log-tail}
 \int_{\R\setminus W_X}K(x-X)|\log|V_\lambda(x+\ii y)||\dd x
 \le T_X:=D_VX^{-3}\log^2X,\qquad y=0,8.
\end{equation}
This is an estimate for the exact logarithm. The comparison zeta log
has a separate exterior bound \(D_ZX^{-4}\), since its absolute
prime-power series is bounded by \(D_0\).

Finally \(\widehat K(\omega)=1/[2\cosh(4\omega)]\), so at
\(\omega=(\log n)/2\) it equals \(1/(n^2+n^{-2})\).
The upper-boundary term less the center term has coefficient
\[
 \frac{n^{-9/2-c/4}}{n^2+n^{-2}}-n^{-5/2-c/4}
       =-\frac{n^{-5/2-c/4}}{1+n^{-4}}.
\]
Absolute convergence permits termwise integration. Equation
\eqref{eq:normalized-green-exact} now has error at most
\(6C_F/L_X+T_X+D_ZX^{-4}\), proving the assertion after enlarging
the onset and taking \(C_{\rm out}=6C_F+D_V+D_Z\).
The prime tail follows from the integral test.
\end{proof}

The constants \(C_{\rm out}\) and \(D_V\) denote the child's
\(C_*\) and \(D_V\); this avoids collision with the earlier
unit-log constant. The onset includes the inherited fixed-buffer
onsets, \(X\ge26\), \(R_X+8\le X/2\),
\(L_XX^{-3/4}\le1\), \(L_X>4C_F\), and
\(L_XX^{-3}\log^2X\le1\). Each gate is eventual after the constants
are fixed. The result holds at every sufficiently large center, but
estimates the outer reduction error, not the sign of \(I-P\).
There is no extra time-zero pole correction in
\eqref{eq:outer-baseline}: the carrier here is harmonic and nonvanishing.

\subsection{Logarithmic transfer without a real-line floor}
On the real window define the two retained arms by
\begin{equation}\label{eq:real-two-arms}
 A_\lambda(x)=\sum_{n\le N_x}w_n(\lambda,x)
                  \ee^{\ii\nu_\lambda(x)\log n},\qquad
 S_\lambda=A_\lambda+\ee^{\ii\phi_\lambda}\overline{A_\lambda},
 \quad\phi_\lambda=2\Theta(\lambda,x).
\end{equation}
The exact field is \(V_\lambda=S_\lambda+e_\lambda\), with
\begin{equation}\label{eq:real-provider-error}
 \sup_{W_X}|e_\lambda|\le C X^{-1/4},\qquad
 C=2^{3/4}(2h_9)+2^{1/4}g_9f_9.
\end{equation}
Indeed the two carriers have equal modulus, the full correction costs
\(g_9f_9x^{-1/4}\), and the provider remainder costs
\(2h_9x^{-3/4}\). Thus the correction is paid in \(e_\lambda\),
not discarded when it is absent from \(S_\lambda\).

\begin{lemma}[One-sided logarithmic transfer]\label{lem:upper-log-transfer}
For \(v,s\in\C\) and \(\delta>0\),
\begin{equation}\label{eq:upper-log-transfer}
 \log|v|\le\tfrac12\log(|s|^2+\delta^2)
       +\tfrac12\log(1+|v-s|^2/\delta^2).
\end{equation}
\end{lemma}
\begin{proof}
Use \(|v|\le|s|+|v-s|\) and Cauchy--Schwarz on
\((|s|,\delta)\) and \((1,|v-s|/\delta)\), then take logarithms.
At \(v=0\) the inequality is understood with \(\log0=-\infty\).
\end{proof}
Combining this lemma with \eqref{eq:real-provider-error}, the exact
tail and Theorem~\ref{thm:outer-baseline}, and choosing
\(\delta=X^{-1/8}\), gives
\begin{align}\label{eq:full-regularized-transfer}
 J_\lambda(X)\le{}&\int_{W_X}\tfrac12
           \log(|S_\lambda|^2+X^{-1/4})\,d\omega_X-P_{c,\lambda}(X)
            +T_X+\frac{C_{\rm out}}{L_X}\nonumber\\
 &+\frac{q_X}{2}\log(1+C^2X^{-1/4}).
\end{align}
The field-transfer cost is \(O(X^{-1/4})\). Neither the exact field
nor the reflected approximant is assumed bounded away from zero.

\subsection{A shifted-zeta comparison retains its unknown divisor}
For comparison with a single ordinary zeta function, let
\(\eta=\pi\lambda/4\) and
\(Z_{c,\lambda}(z)=\zeta((1-\ii z)/2+c/4-\ii\pi\lambda/8)\).
The exact time-zero normalization gives
\[
 M((1-\ii z)/2+c/4-\ii\pi\lambda/8)Z_{c,\lambda}(z)
       =H_0(z+\eta+\ii c/2).
\]
In the strip the zeta factor has a simple pole
\(p=-\eta+\ii(1-c/2)\); its gamma completion has precisely the
cancelling simple zero there. If \(x_j(0)+\ii y_j(0)\) are the
time-zero zeros, with multiplicity, set
\[
 \mathcal Z_{c,\lambda}(X)=\sum_{y_j(0)>c/2}
       G(X+\eta-x_j(0),y_j(0)-c/2),\qquad
 \Delta=I_\lambda-\int_\R K(x-X)\log|Z_{c,\lambda}(x)|\dd x.
\]
The complete-product proof of Theorem~\ref{thm:green} applies to this
fixed affine translate: local distance logarithms remain integrable,
the count bound is unchanged, and the boundary tails remain summable.
The pole enters with negative sign. On the two outer lines the zeta
combination is exactly \(-P_{c,\lambda}\); therefore
\begin{equation}\label{eq:shifted-zeta-divisor}
 I_\lambda-P_{c,\lambda}
    =\Delta+\mathcal Z_{c,\lambda}-G(X+\eta,1-c/2).
\end{equation}
This is B02's meromorphic identity \cite{GNPRSComparisons}. The
difference \(\Delta\) is a difference of integrable logarithms, not
an assumed holomorphic logarithm of a zero-free quotient. A useful
upper bound still requires a coupled estimate for
\(\Delta+\mathcal Z_{c,\lambda}\). The identity alone supplies none;
the isolated pole term decays exponentially with \(X\).
